% checkpoint.tex
% Relatório de checkpoint da iniciação científica
% trajectory-prediction / drift_analise
%
% Compilar com:
%   pdflatex checkpoint.tex
% (não usa bibliografia externa; referências em \bibitem)

\documentclass[11pt,a4paper]{article}

\usepackage[utf8]{inputenc}
\usepackage[T1]{fontenc}
\usepackage[brazil]{babel}
\usepackage{lmodern}
\usepackage{microtype}
\usepackage{geometry}
\geometry{margin=2.5cm}
\usepackage{amsmath,amssymb}
\usepackage{booktabs}
\usepackage{enumitem}
\usepackage{hyperref}
\hypersetup{colorlinks=true, linkcolor=black, urlcolor=blue, citecolor=black}
\usepackage{xcolor}
\usepackage{listings}
\lstset{
    basicstyle=\ttfamily\small,
    breaklines=true,
    columns=fullflexible,
    frame=single,
    framesep=4pt,
    xleftmargin=4pt,
    aboveskip=6pt,
    belowskip=6pt,
}
\usepackage{titlesec}
\titlespacing*{\section}{0pt}{1.4em}{0.6em}
\titlespacing*{\subsection}{0pt}{1.0em}{0.4em}

\title{\vspace{-1.5cm}Checkpoint --- An\'alise de Data Drift em Seq2Seq para
       Predi\c{c}\~ao de Trajet\'orias no RoboCup SSL}
\author{Iniciação Científica --- Relatório de progresso (Meses 1--4)}
\date{\today}

\begin{document}
\maketitle
\vspace{-0.6cm}

% ============================================================================
\section*{Resumo}
% ============================================================================

Este checkpoint documenta as mudan\c{c}as realizadas no projeto
\texttt{trajectory-prediction} entre a implementa\c{c}\~ao do M\^es 1 e o
estado atual da pesquisa. O objetivo do M\^es 1 foi estabelecer
infraestrutura granular (por jogo, por janela e por trajet\'oria) para
substituir a granularidade anual do notebook original. O M\^es 2
introduziu detectores formais de drift (ADWIN, Page-Hinkley, KSWIN, Pelt)
para trocar a infer\^encia visual por testes estat\'isticos. Em seguida,
uma rodada cr\'itica identificou 19 pontos fr\'ageis no pipeline e
implementou as corre\c{c}\~oes priorit\'arias: ordem temporal das
trajet\'orias, recalibra\c{c}\~ao de detectores, decomposi\c{c}\~ao
covariate vs concept drift, null detector via permuta\c{c}\~ao e curva de
lat\^encia. O notebook unificado tem hoje 55 c\'elulas e produz 39 figuras
em PT/EN al\'em de 12 CSVs de m\'etricas, todos versionados em
\texttt{Relas/results/\{mes1,mes2\}/}.

O M\^es~3 implementou adapta\c{c}\~ao ao drift por dois caminhos:
(i)~importance weighting via LSIF, que decompõe o excesso de ADE em
componentes covariate vs.\ concept drift; e (ii)~retreino seletivo
nos breakpoints detectados pelo Pelt. Os resultados indicam que 59--72\,\%
do excesso \'e covariate (feature dominante: \texttt{accel\_p99}), mas o
fine-tuning induz \emph{catastrophic interference} severo
(\texttt{cat\_ratio}\,=\,0{,}82 j\'a na 1\textordfeminine{} \'epoca). A
valida\c{c}\~ao por divis\~ao confirmou que o drift \'e real e
intra-divis\~ao (Divis\~ao\,B ausente nos dados). O M\^es~4 consolidou
os resultados, aplicou a \'arvore de decis\~ao de cen\'ario e produziu
este relat\'orio final com tese central (Cen\'ario~3 descritivo).

% ============================================================================
\section{Linha do tempo}
% ============================================================================

\begin{enumerate}[leftmargin=2em,itemsep=2pt]
    \item \textbf{M\^es 1 --- Infraestrutura granular.} Enriquecimento do
          \texttt{dataset.csv} com contexto cinem\'atico por jogo, expans\~ao
          para granularidade de janela e amostra de erros por trajet\'oria.
    \item \textbf{M\^es 2 --- Detectores formais.} ADWIN, Page-Hinkley,
          KSWIN, Pelt aplicados sobre os streams definidos no M\^es 1.
    \item \textbf{Itera\c{c}\~ao cr\'itica.} An\'alise de 19 pontos
          fracos; reescrita do notebook unificado e do
          \texttt{plot\_covariate\_shift.py}; cria\c{c}\~ao de
          \texttt{compute\_trajectory\_errors.py}.
    \item \textbf{Estado atual.} Pipeline documentado em
          \texttt{Relas/PIPELINE.md}, dependências fixadas em
          \texttt{requirements.txt}, modelo configurado para target
          \texttt{dv} (delta-velocity), divis\~ao A/B mapeada via
          \texttt{division\_map.csv}. Pronto para o M\^es 3.
    \item \textbf{M\^es 3 --- Adapta\c{c}\~ao ao drift.} Importance
          weighting (LSIF) com decomposi\c{c}\~ao por feature;
          retreino seletivo no breakpoint Pelt com early-stop por
          \texttt{cat\_ratio}; valida\c{c}\~ao por divis\~ao SSL
          (2021 mapeado como Both; Divis\~ao\,B ausente).
    \item \textbf{M\^es 4 --- Consolida\c{c}\~ao e reda\c{c}\~ao.}
          Re-execu\c{c}\~ao com $n_{\text{por ano}}=6$; decis\~ao da
          tese central (Cen\'ario~3 descritivo); atualiza\c{c}\~ao
          deste documento e do \texttt{main\_pt.tex}.
\end{enumerate}

% ============================================================================
\section{M\^es 1 --- Infraestrutura granular}
% ============================================================================

\subsection{Entregas}

\begin{itemize}[leftmargin=1.4em,itemsep=2pt]
    \item \texttt{plot\_covariate\_shift.py} estendido para gerar dois
          arquivos novos:
        \begin{itemize}[leftmargin=1.4em,itemsep=1pt]
            \item \texttt{per\_game\_context.csv} --- uma linha por jogo
                  com \texttt{match\_id}, \texttt{log\_file}, ano, divis\~ao,
                  \texttt{avg\_robots\_per\_stoppage}, n\'umero de
                  trajet\'orias, dura\c{c}\~ao estimada e percentis
                  cinem\'aticos (\texttt{speed\_p\{50,90,99\}},
                  \texttt{accel\_p\{50,90,99\}}, \texttt{turn\_p\{50,90\}}).
            \item \texttt{per\_window\_drift.csv} --- KS e Wasserstein por
                  janela contra o \emph{baseline} de treino
                  (\texttt{proc\_set\_1+2}). Tamanho de janela
                  \texttt{N\_TRAJS\_PER\_WINDOW = 30} (calibrado abaixo).
        \end{itemize}
    \item \texttt{compute\_trajectory\_errors.py} (script novo) --- carrega
          o Seq2Seq treinado e o \emph{Kalman} determin\'istico, roda numa
          amostra de 3 jogos por ano (\texttt{seed = 42}) e produz
          \texttt{trajectory\_errors\_sample.parquet} com ADE/FDE por
          janela.
    \item Notebook \texttt{drift\_analise.ipynb} (Se\c{c}\~oes 1--5) com
          descritivas na nova granularidade: ADE por jogo com IC
          \emph{bootstrap}, ADE vs features cinem\'aticas, heatmap
          Spearman, tabelas de resumo e checagens de sanidade.
    \item Sa\'idas em PT e EN salvas em \texttt{Relas/results/mes1/}: 19
          figuras PDF (par PT/EN) + tr\^es CSVs
          (\texttt{summary\_granularity.csv}, \texttt{tail\_stats.csv},
          \texttt{window\_intra\_year\_var.csv}).
\end{itemize}

\subsection{Decis\~oes de design}

\begin{itemize}[leftmargin=1.4em,itemsep=2pt]
    \item \emph{Granularidade de janela}: a especifica\c{c}\~ao inicial
          fixava $N=500$ trajet\'orias, o que produzia uma janela por jogo
          em quase todos os casos --- equivalente \`a granularidade anual
          do \texttt{ks\_wd\_per\_dataset.csv} j\'a existente. Reduzimos
          para $N=30$ ap\'os a itera\c{c}\~ao cr\'itica, obtendo mediana
          de 6 janelas por jogo (Tabela~\ref{tab:granularity}).
    \item \emph{Ordem temporal}: a vers\~ao original de
          \texttt{iter\_robot\_trajs} concatenava todas as trajet\'orias
          do time \texttt{blue} antes do \texttt{yellow}, sem garantia
          de ordem cronol\'ogica. Foi reescrita para ordenar por
          \texttt{time\_c[0]} de cada trajet\'oria.
    \item \emph{Bilinguismo}: todos os plots s\~ao gerados em PT e EN via
          \texttt{T(pt, en, lang)}, salvando o par
          \texttt{<nome>\_\{pt,en\}.pdf} em \texttt{Relas/results/mes1/}.
\end{itemize}

\begin{table}[h]
\centering
\small
\begin{tabular}{rrrr}
\toprule
\textbf{Ano} & \textbf{N\'umero de jogos} & \textbf{Trajet\'orias totais} & \textbf{Janelas totais} \\
\midrule
2019 & 6 & 1\,171 & 42 \\
2021 & 6 & 2\,206 & 77 \\
2022 & 6 &   807 & 30 \\
2023 & 6 &   903 & 34 \\
2024 & 6 &   886 & 31 \\
2025 & 6 & 1\,055 & 39 \\
\midrule
\textbf{Total} & \textbf{36} & \textbf{7\,028} & \textbf{253} \\
\bottomrule
\end{tabular}
\caption{Granularidade obtida com \texttt{N\_TRAJS\_PER\_WINDOW = 30}.
         Fonte: \texttt{Relas/results/mes1/summary\_granularity.csv}.}
\label{tab:granularity}
\end{table}

% ============================================================================
\section{M\^es 2 --- Detectores formais de drift}
% ============================================================================

\subsection{Entregas}

\begin{itemize}[leftmargin=1.4em,itemsep=2pt]
    \item \textbf{Se\c{c}\~ao 6 --- Concept drift online.} ADWIN
          (\texttt{delta = 0.002}) e Page-Hinkley
          (\texttt{threshold = 5.0 mm}, recalibrado pela ordem de grandeza
          do ADE) sobre a s\'erie de \texttt{ade\_traj} do Seq2Seq
          $30{\to}15$, ordenada por \texttt{(year, match\_id, traj\_id)}.
          O mesmo procedimento \'e aplicado ao Kalman determin\'istico para
          permitir leitura cruzada Seq2Seq vs Kalman --- drift no Kalman
          sinaliza mudan\c{c}a no comportamento do alvo, drift apenas no
          Seq2Seq sinaliza miscalibra\c{c}\~ao do modelo.
    \item \textbf{Se\c{c}\~ao 7 --- Covariate shift online.} KSWIN com
          \texttt{window\_size = 30}, \texttt{stat\_size = 10},
          \texttt{alpha = 0.005} sobre \texttt{ks\_speed\_window} e
          \texttt{ks\_accel\_window}. Par\^ametros calibrados para o stream
          de 253 janelas; com $N=500$ original n\~ao havia amostras
          suficientes para o detector disparar.
    \item \textbf{Se\c{c}\~ao 8 --- Benchmark offline com Pelt.}
          \texttt{ruptures.Pelt} com custo RBF sobre o ADE m\'edio por
          jogo (s\'erie de 30 pontos). \emph{Elbow plot} variando a
          penalidade $pen \in \{1, 2, 3, 5, 7, 10, 15, 20, 30, 50, 75, 100, 150\}$
          guia a escolha autom\'atica de \texttt{PEN\_CHOSEN}.
    \item \textbf{Se\c{c}\~ao 9 --- S\'intese.} Fun\c{c}\~ao
          \texttt{map\_idx\_to\_game} converte \'indices de cada stream
          em \texttt{(year, match\_id, log\_file)}. Tabela
          \texttt{alarms\_consolidated.csv} re\'une os alarmes de todos
          os detectores para discuss\~ao.
    \item Sa\'idas em \texttt{Relas/results/mes2/}.
\end{itemize}

% ============================================================================
\section{Itera\c{c}\~ao cr\'itica e melhorias incorporadas}
% ============================================================================

A revis\~ao do estado p\'os-M\^es 2 identificou 19 pontos fr\'ageis,
agrupados em tr\^es categorias: integridade dos dados, calibra\c{c}\~ao
estat\'istica e an\'alises faltantes. As melhorias priorit\'arias foram
implementadas:

\subsection{Integridade dos dados}

\begin{itemize}[leftmargin=1.4em,itemsep=2pt]
    \item \emph{Ordem temporal das trajet\'orias.} \texttt{iter\_robot\_trajs}
          agora intercala \texttt{blue} e \texttt{yellow} ordenando por
          \texttt{time\_c[0]} (par\^ametro \texttt{chrono = True}).
    \item \emph{Granularidade de janela.} $N=500 \to N=30$.
    \item \emph{Coluna \texttt{n\_robots}.} Renomeada para
          \texttt{avg\_robots\_per\_stoppage}, refletindo que \'e uma
          \emph{proxy} (raz\~ao trajet\'orias / \texttt{stop\_id}\'s
          \'unicos), n\~ao o n\'umero real de rob\^os simult\^aneos em
          campo.
    \item \emph{Modo de predi\c{c}\~ao.} \texttt{compute\_trajectory\_errors.py}
          fixa \texttt{target = "dv"} (delta-velocity) para alinhar com
          os pesos treinados; o \emph{Kalman} de refer\^encia passa a ser
          $\hat{\Delta v} = 0$ (velocidade constante implica incremento
          nulo).
\end{itemize}

\subsection{Calibra\c{c}\~ao estat\'istica}

\begin{itemize}[leftmargin=1.4em,itemsep=2pt]
    \item \emph{Page-Hinkley.} \texttt{threshold} ajustado de 50 para
          5\,mm. Para ADE em mil\'imetros com mediana
          \mbox{$\sim 10$--$50$}\,mm, threshold $= 50$ era praticamente
          imposs\'ivel de disparar. O novo valor torna o detector vis\'ivel
          ao ru\'ido real do dom\'inio.
    \item \emph{KSWIN.} \texttt{window\_size} ajustado para 30 e
          \texttt{stat\_size} para 10, compat\'iveis com o tamanho efetivo
          do stream de janelas.
    \item \emph{Pelt.} Penalidade escolhida via \emph{elbow plot} ao
          inv\'es de chute fixo.
\end{itemize}

\subsection{An\'alises adicionadas}

\begin{itemize}[leftmargin=1.4em,itemsep=2pt]
    \item \textbf{Se\c{c}\~ao 10.1 --- Null detector.} 50 permuta\c{c}\~oes
          do stream e contagem de alarmes por permuta\c{c}\~ao. Reporta
          $p$-valor emp\'irico $P(\text{null} \geq \text{observado})$.
          Sem esta etapa, qualquer alarme observado seria interpretado
          como evid\^encia, mesmo se compat\'ivel com ru\'ido temporal.
    \item \textbf{Se\c{c}\~ao 10.2 --- Lat\^encia de detec\c{c}\~ao.} Stream
          \emph{bootstrap} de 2019 com shift sint\'etico injetado em
          $t^* = 200$, varrendo
          $\Delta \in \{0, 1, 2, 5, 10, 20, 50\}$\,mm. Para cada par
          (detector, $\Delta$) reporta \texttt{lat\_p10}, \texttt{lat\_p50},
          \texttt{lat\_p90}. M\'etrica essencial para o M\^es 3, que
          decidir\'a quando acionar retreino.
    \item \textbf{Se\c{c}\~ao 11 --- Decomposi\c{c}\~ao covariate vs
          concept.} Resposta para a pergunta central da IC:
          \begin{equation}
              \underbrace{\text{ADE}_y - \text{ADE}_{2019}}_{\text{excess total}}
              = \underbrace{\text{ADE}^{\text{compat}}_y - \text{ADE}_{2019}}_{\text{concept}}
              + \underbrace{(\text{ADE}_y - \text{ADE}^{\text{compat}}_y)}_{\text{covariate}}
              \notag
          \end{equation}
          onde $\text{ADE}^{\text{compat}}_y$ \'e o ADE m\'edio em jogos
          do ano $y$ com $\texttt{ks\_speed} \leq K^*$ e $K^*$ \'e a
          mediana de \texttt{ks\_speed} em 2019. Sa\'ida em
          \texttt{drift\_decomposition.csv}.
    \item \textbf{Comparativo Seq2Seq vs Kalman.} ADWIN/Page-Hinkley
          rodam em ambos os streams. Diferen\c{c}a entre os dois isola
          miscalibra\c{c}\~ao do modelo do drift do alvo.
\end{itemize}

% ============================================================================
\section{Estado atual do pipeline}
% ============================================================================

A Tabela~\ref{tab:pipeline} resume a cadeia de produ\c{c}\~ao. O fluxo
completo est\'a documentado em \texttt{Relas/PIPELINE.md} com
\emph{data dictionary} de todos os artefatos.

\begin{table}[h]
\centering
\small
\begin{tabular}{lll}
\toprule
\textbf{Origem} & \textbf{Script} & \textbf{Destino} \\
\midrule
\texttt{dataset/proc\_set\_*.pkl} & \texttt{plot\_covariate\_shift.py} &
   \texttt{covariate\_shift\_out/per\_game\_context.csv} \\
                                  &                                    &
   \texttt{covariate\_shift\_out/per\_window\_drift.csv} \\
                                  &                                    &
   \texttt{covariate\_shift\_out/ks\_wd\_per\_dataset.csv} \\
\midrule
\texttt{robot\_30\_15\_t.weights.h5} & \texttt{compute\_trajectory\_errors.py} &
   \texttt{covariate\_shift\_out/trajectory\_errors\_sample.parquet} \\
\texttt{normalization\_stats\_*.pkl} &                                  & \\
\midrule
artefatos acima                  & \texttt{drift\_analise.ipynb} &
   \texttt{Relas/results/mes1/*} \\
                                  &                                    &
   \texttt{Relas/results/mes2/*} \\
                                  &                                    &
   \texttt{dataset\_enriched.csv} \\
\bottomrule
\end{tabular}
\caption{Cadeia de produ\c{c}\~ao do pipeline.}
\label{tab:pipeline}
\end{table}

\subsection{Como reproduzir}

\begin{lstlisting}[language=bash]
pip install -r requirements.txt
python drift_analise/plot_covariate_shift.py
python model_analise/compute_trajectory_errors.py --n_per_year 3 --seed 42
jupyter nbconvert --to notebook --execute drift_analise/drift_analise.ipynb \
    --output drift_analise.executed.ipynb --output-dir drift_analise \
    --ExecutePreprocessor.timeout=600
\end{lstlisting}

\subsection{Dependências novas}

\texttt{river == 0.21.2} (detectores online) e
\texttt{ruptures == 1.1.9} (detec\c{c}\~ao offline) foram adicionadas
ao \texttt{requirements.txt}.

% ============================================================================
\section{An\'alise da pesquisa atual}
% ============================================================================

\subsection{O que o pipeline j\'a permite afirmar}

\begin{enumerate}[leftmargin=1.6em,itemsep=2pt]
    \item Existe degrada\c{c}\~ao de ADE entre 2019 (treino) e os anos
          subsequentes, vis\'ivel tanto em granularidade anual
          (notebook original) quanto por jogo (Se\c{c}\~ao 4 do notebook
          atual).
    \item A degrada\c{c}\~ao tem cauda mais pesada do que m\'edia: a
          raz\~ao \texttt{p99/p50} do ADE cresce mais r\'apido do que a
          m\'edia (Figura \texttt{3c\_p99\_p50\_ratio}). O modelo n\~ao
          erra ``um pouco mais em tudo''; ele falha cataclismicamente em
          subconjuntos espec\'ificos.
    \item As features cinem\'aticas \texttt{ks\_speed} e
          \texttt{ks\_accel} variam intra-ano (Figuras
          \texttt{2b\_window\_box\_*}), o que justifica granularidade
          fina e n\~ao s\'o m\'edia anual.
    \item ADE Seq2Seq por jogo correlaciona positivamente com
          \texttt{speed\_p90} e \texttt{accel\_p90} do jogo
          (Figuras \texttt{4b\_*}), sugerindo que jogos com cauda
          cinem\'atica mais agressiva s\~ao sistematicamente piores
          para o modelo.
\end{enumerate}

\subsection{O que ainda \'e hip\'otese}

\begin{enumerate}[leftmargin=1.6em,itemsep=2pt]
    \item \emph{Quem domina: covariate ou concept drift?}
          A Se\c{c}\~ao 11 do notebook produz uma estimativa, mas o valor
          definitivo depende de rodar
          \texttt{compute\_trajectory\_errors.py} com os pesos reais (e
          n\~ao com o \emph{stub} sint\'etico usado durante o
          desenvolvimento) e do tamanho da amostra (3 jogos por ano \'e
          frugal demais). \textbf{Pr\'oxima a\c{c}\~ao}: aumentar para
          $\geq 5$ jogos por ano e refazer a Tabela 11.
    \item \emph{Os alarmes coincidem com mudan\c{c}as conhecidas da
          liga?} A discuss\~ao da Se\c{c}\~ao 12 lista hip\'oteses
          (hiato pandêmico 2019$\to$2022, evolu\c{c}\~ao de hardware,
          mudan\c{c}as de regras), mas confirma\c{c}\~ao requer cruzar
          com o \emph{rule-book changelog} oficial da SSL e os
          \emph{Champion Team Description Papers} dos anos
          relevantes~\cite{rcssl_rules}.
    \item \emph{Os detectores t\^em sinal acima do null?} A Se\c{c}\~ao
          10.1 responde quantitativamente, mas ainda n\~ao foi rodada
          com o stream real --- depende novamente de
          \texttt{trajectory\_errors\_sample.parquet} produzido pelo
          modelo verdadeiro.
\end{enumerate}

\subsection{Limita\c{c}\~oes honestas}

\begin{itemize}[leftmargin=1.4em,itemsep=2pt]
    \item \textbf{Tamanho amostral.} 3 jogos por ano $\times$ 7 anos = 21
          jogos efetivos (excluindo 2020). Para detectores que se beneficiam
          de \emph{streams} grandes (ADWIN, Page-Hinkley), isso \'e o
          \'unico recurso para concept drift; para covariate shift,
          \texttt{df\_win} d\'a 253 janelas, o que \'e mais confort\'avel.
    \item \textbf{Saltos artificiais entre anos.} O stream concatena anos
          em sequ\^encia (2019 $\to$ 2021 $\to$ 2022 $\to \ldots$). Os
          detectores online v\~ao reportar drift na \emph{fronteira} de
          ano, mas isso \'e parcialmente artefato da concatena\c{c}\~ao
          n\~ao-cont\'inua. As fronteiras est\~ao destacadas nos plots
          como linhas pontilhadas para leitura honesta.
    \item \textbf{Calibra\c{c}\~ao geral\'istica.} Todos os detectores
          usam defaults da literatura, recalibrados quando vis\'ivelmente
          impertinentes (Page-Hinkley com \texttt{threshold = 50} era
          cego). Calibra\c{c}\~ao espec\'ifica do dom\'inio SSL (\emph{e.g.}
          via valida\c{c}\~ao cruzada em \emph{slices} sint\'eticos)
          \'e tarefa do M\^es 3.
    \item \textbf{Ground truth do drift.} N\~ao temos r\'otulos de
          ``drift sim/n\~ao'' para cada janela. Toda a valida\c{c}\~ao
          \'e indireta: null permutacional (Se\c{c}\~ao 10.1) e injection
          de shift sint\'etico (Se\c{c}\~ao 10.2).
\end{itemize}

% ============================================================================
\section{M\^es 3 --- Adapta\c{c}\~ao ao drift}
% ============================================================================

\subsection{Importance weighting (LSIF)}

O objetivo do M\^es~3 foi quantificar quanto do excesso de ADE \'e
explic\'avel por \emph{covariate shift} (mudan\c{c}a na distribui\c{c}\~ao
das entradas) em contraste com \emph{concept drift} (mudan\c{c}a em
$p(y|x)$). O m\'etodo adotado foi LSIF (\emph{Least-Squares Importance
Fitting}), que estima a raz\~ao de densidades
$\hat{w}(x) = p_{\text{alvo}}(x) / p_{2019}(x)$ diretamente,
sem estimar as densidades individuais~\cite{kliep,sugiyama_book}.
As \emph{features} s\~ao os oito percentis cinem\'aticos por jogo
(\texttt{speed/accel/turn} $\times$ \{mean, p90, p99\} menos
\texttt{turn\_p99}).

O ADE ponderado e a \emph{recovery} s\~ao definidos por:
\begin{equation}
\mathrm{ADE}_{\mathrm{IW}}(y) =
    \frac{\displaystyle\sum_{i} \hat{w}(x_i)\,e_i}
         {\displaystyle\sum_{i} \hat{w}(x_i)}, \qquad
\mathrm{recovery}(y) =
    \frac{\mathrm{ADE}_y - \mathrm{ADE}_{\mathrm{IW}}(y)}
         {\mathrm{ADE}_y - \mathrm{ADE}_{2019}} \times 100\,\%.
\label{eq:iw}
\end{equation}
Se o drift for puramente covariate, $\mathrm{ADE}_{\mathrm{IW}} \approx
\mathrm{ADE}_{2019}$ e recovery $\approx 100\,\%$.
O ESS ratio $\bigl[(\sum_i \hat{w}_i)^2 / (n \sum_i \hat{w}_i^2)\bigr]$
mede a estabilidade do estimador IW.

\begin{table}[h]
\centering
\small
\begin{tabular}{r r r r r r}
\toprule
\textbf{Ano} & \textbf{excess (mm)} & \textbf{ADE\textsubscript{IW} (mm)}
             & \textbf{ESS} & \textbf{recovery (\%)} & \textbf{IC\,95\,\%} \\
\midrule
2021 & $+3{,}05$ & 5{,}70 & 0{,}58 & 61{,}6 & [60{,}1;\;63{,}1] \\
2022 & $+0{,}23$ & 4{,}62 & 0{,}73 & 59{,}0 & [46{,}5;\;77{,}9] \\
2023 & $+0{,}69$ & 4{,}72 & 0{,}72 & 71{,}8 & [66{,}9;\;77{,}4] \\
2024 & $-0{,}03$ & ---     & 0{,}79 & ---    & --- \\
2025 & $+0{,}73$ & 4{,}75 & 0{,}59 & 69{,}3 & [64{,}4;\;75{,}1] \\
\bottomrule
\end{tabular}
\caption{Decomposi\c{c}\~ao IW ($n_{\text{source}} = n_{\text{target}}
         = 48\,000$; \texttt{--n\_per\_year~6}). ESS\,$>$\,0{,}5 indica
         estimador est\'avel. Ano 2024 exclu\'ido (sem excesso,
         ADE$_{2024}$ $<$ ADE$_{2019}$). IC\,95\,\% via bootstrap
         com 200 amostras.}
\label{tab:iw_decomp}
\end{table}

\textbf{Decomposi\c{c}\~ao por feature.} Para identificar quais dimens\~oes
dominam o covariate shift, o IW foi re-estimado usando cada feature
individualmente (\texttt{--per\_feature}). Ranking m\'edio de
recovery\_pct entre 2021, 2023 e 2025:
(1)~\texttt{accel\_p99} $\approx 54\,\%$,
(2)~\texttt{accel\_p90} $\approx 52\,\%$,
(3)~\texttt{accel\_mean} $\approx 47\,\%$,
(4--6)~features de velocidade ($20$--$40\,\%$),
(7--8)~\texttt{turn\_\{mean,p90\}} ($<5\,\%$).
O ranking estabiliza entre anos: acelera\c{c}\~ao de cauda
(\texttt{accel\_p99}) \'e a dimens\~ao de covariate shift mais
relevante para o modelo Seq2Seq.

\subsection{Retreino seletivo}

O Pelt (pen\,=\,5) detectou um breakpoint inter-ano: transi\c{c}\~ao
2019+2021 $\to$ 2022--2025 (\texttt{antes\_2022}). O pipeline de
fine-tuning: (i)~congela o encoder e as primeiras camadas do decoder;
(ii)~descongela as 2 camadas finais; (iii)~Adam lr\,=\,$10^{-5}$,
at\'e 20 \'epocas com early-stop em \texttt{cat\_ratio}\,$>$\,0{,}5.
O \texttt{cat\_ratio} mede a raz\~ao trajet\'orias degradadas /
melhoradas, capturando \emph{catastrophic interference}.

\begin{table}[h]
\centering
\small
\begin{tabular}{l c c c r r r}
\toprule
\textbf{Breakpoint} & \textbf{n\_unfreeze} & \textbf{lr}
    & \textbf{\'epocas} & \textbf{$\Delta$ADE (mm)}
    & \textbf{recovery (\%)} & \textbf{cat\_ratio} \\
\midrule
\texttt{antes\_2022} & 2 & $10^{-5}$ & 1 (stop) & $-0{,}13$ & 32{,}7 & 0{,}82 \\
\bottomrule
\end{tabular}
\caption{Resultados do retreino seletivo. Treino interrompido na 1\textordfeminine{} \'epoca
         pois \texttt{cat\_ratio}\,=\,0{,}82\,$>$\,0{,}5 (limiar).}
\label{tab:retrain}
\end{table}

O fine-tuning reduziu o ADE m\'edio em apenas $0{,}13$\,mm ($-2{,}7\,\%$),
recovery\,=\,32{,}7\,\%. Para cada trajet\'oria melhorada, 0{,}82 foram
pioradas, indicando que a adapta\c{c}\~ao da m\'edia vem \`a custa de
degrada\c{c}\~ao severa em trajet\'orias individuais.

% ============================================================================
\section{Valida\c{c}\~ao por divis\~ao}
% ============================================================================

O script \texttt{division\_validation.py} mapeou cada
\texttt{proc\_set} \`a sua divis\~ao via \texttt{division\_map.csv}.
Resultado: todos os jogos p\'os-2021 pertencem \`a Divis\~ao\,A;
a Divis\~ao\,B n\~ao est\'a representada. O ano 2021 foi classificado
como \emph{Both} (formato h\'ibrido p\'os-pandemia).

\begin{table}[h]
\centering
\small
\begin{tabular}{l r r r r r}
\toprule
\textbf{Divis\~ao} & \textbf{Ano}
    & \textbf{ADE\textsubscript{2019} (mm)} & \textbf{ADE\textsubscript{y} (mm)}
    & \textbf{excess (mm)} & \textbf{jogos} \\
\midrule
Both & 2021 & ---  & 7{,}58 & ---       & 6 \\
A    & 2022 & 4{,}53 & 4{,}75 & $+0{,}23$ & 6 \\
A    & 2023 & 4{,}53 & 5{,}21 & $+0{,}69$ & 6 \\
A    & 2024 & 4{,}53 & 4{,}50 & $-0{,}03$ & 6 \\
A    & 2025 & 4{,}53 & 5{,}26 & $+0{,}73$ & 6 \\
\bottomrule
\end{tabular}
\caption{Excesso de ADE por divis\~ao. Baseline: ADE$_{2019}$\,=\,4{,}53\,mm
         (Divis\~ao\,A, 6 jogos, $n=6$).}
\label{tab:division}
\end{table}

Como todos os dados p\'os-2021 s\~ao da Divis\~ao\,A, o excesso m\'edio
($\approx 0{,}40$\,mm) n\~ao pode ser atribu\'ido \`a mudan\c{c}a de
divis\~ao: o drift \'e real e intra-divis\~ao. A aus\^encia da
Divis\~ao\,B impede comparar regimes competitivos distintos.

% ============================================================================
\section{An\'alise cr\'itica e tese central}
% ============================================================================

\subsection{\'Arvore de decis\~ao}

A escolha do cen\'ario seguiu a \'arvore definida no plano do M\^es~4:

\begin{itemize}[leftmargin=1.4em,itemsep=2pt]
    \item \textbf{Cen\'ario~1} (``drift \'e confundimento de divis\~ao''):
          requer $\Delta_A < 0{,}3$\,mm \emph{e} $\Delta_B > 1{,}5$\,mm.
          \textbf{N\~ao se aplica:} $\Delta_A \approx 0{,}40$\,mm e
          Divis\~ao\,B ausente.
    \item \textbf{Cen\'ario~2} (``detectar \'e f\'acil, adaptar n\~ao''):
          requer $\Delta_A > 0{,}5$\,mm \emph{e} recovery m\'ediana
          $< 50\,\%$. \textbf{N\~ao se aplica:} $\Delta_A = 0{,}40$\,mm
          e recovery m\'ediana $\approx 63\,\%$ (acima do limiar).
    \item \textbf{Cen\'ario~3} (descritivo sem afirma\c{c}\~ao forte):
          \textbf{escolhido por exclus\~ao.}
\end{itemize}

\subsection{Tese central --- Cen\'ario~3}

O drift no Seq2Seq para trajet\'orias SSL \'e \textbf{real},
\textbf{intra-divis\~ao} e \textbf{predominantemente covariate}:
entre 59 e 72\,\% do excesso de ADE \'e explicado por mudan\c{c}a na
distribui\c{c}\~ao das features cinem\'aticas (especialmente
\texttt{accel\_p99}). Contudo, a adapta\c{c}\~ao por fine-tuning \'e
severamente prejudicada por \emph{catastrophic interference}.

\textbf{Suporte num\'erico:}
\begin{enumerate}[leftmargin=1.6em,itemsep=2pt]
    \item $\Delta_A \approx 0{,}40$\,mm positivo e persistente em
          2022--2025 (exceto 2024): drift real, n\~ao pontual.
    \item Recovery m\'ediana $\approx 63\,\%$ (IC\,95: 59--72\,\%):
          covariate shift explica a maior parte do excesso, mas
          $\approx 37\,\%$ permanece ap\'os repondera\c{c}\~ao,
          sugerindo algum concept drift residual.
    \item \texttt{accel\_p99} \'e a feature mais informativa do
          covariate shift (recovery individual\,$>$\,50\,\% em todos
          os anos com excesso, ranking est\'avel entre anos).
    \item Fine-tuning com 2 camadas descongeladas e lr\,=\,$10^{-5}$:
          \texttt{cat\_ratio}\,=\,0{,}82, early-stop na 1\textordfeminine{} \'epoca;
          o trade-off m\'edia--cauda \'e desfavor\'avel.
\end{enumerate}

\textbf{Limita\c{c}\~oes relevantes:} Divis\~ao\,B ausente; um \'unico
breakpoint retreinado (Pelt detectou apenas a transic\~ao 2019+2021
$\to$ 2022); fine-tuning testado apenas no decoder. Estas limita\c{c}\~oes
s\~ao candidatas naturais para trabalho futuro.

% ============================================================================
\section*{Lista de artefatos atuais}
% ============================================================================

\begin{itemize}[leftmargin=1.4em,itemsep=1pt,parsep=0pt]
    \item \textbf{Scripts}: \texttt{plot\_covariate\_shift.py},
          \texttt{compute\_trajectory\_errors.py},
          \texttt{compute\_importance\_weights.py},
          \texttt{plot\_iw\_results.py},
          \texttt{retrain\_at\_breakpoints.py},
          \texttt{division\_validation.py}.
    \item \textbf{Notebook}: \texttt{drift\_analise.ipynb} ($\geq$60 c\'elulas).
    \item \textbf{Documenta\c{c}\~ao}: \texttt{Relas/PIPELINE.md},
          \texttt{Relas/CODEX\_CHECKLIST.md}, este arquivo
          (\texttt{Relas/checkpoint.tex}).
    \item \textbf{Configura\c{c}\~ao auxiliar}: \texttt{division\_map.csv},
          \texttt{requirements.txt} (com \texttt{river}, \texttt{ruptures},
          \texttt{scikit-learn}).
    \item \textbf{Sa\'idas M\^es 1}: 19 pares de PDF (PT/EN) e 3 CSVs em
          \texttt{Relas/results/mes1/}.
    \item \textbf{Sa\'idas M\^es 2}: 12 pares de PDF e 6 CSVs em
          \texttt{Relas/results/mes2/}.
    \item \textbf{Sa\'idas M\^es 3}: \texttt{iw\_decomposition.csv},
          \texttt{iw\_per\_feature.csv}, \texttt{retrain\_results.csv},
          \texttt{retrain\_epoch\_log\_bp0.csv},
          \texttt{by\_division\_decomposition.csv},
          \texttt{sample\_size\_sensitivity.csv}, 3 PDF de IW e 1 PDF de
          divis\~ao em \texttt{Relas/results/mes3/}.
\end{itemize}

% ============================================================================
\begin{thebibliography}{9}
% ============================================================================

\bibitem{adwin}
A.~Bifet and R.~Gavald\`a.
\newblock Learning from time-changing data with adaptive windowing.
\newblock In \emph{SIAM International Conference on Data Mining (SDM)}, 2007.

\bibitem{page_hinkley}
E.~S.~Page.
\newblock Continuous inspection schemes.
\newblock \emph{Biometrika}, 41(1/2):100--115, 1954.

\bibitem{kswin}
C.~Raab, M.~Heusinger, F.-M.~Schleif.
\newblock Reactive Soft Prototype Computing for Concept Drift Streams.
\newblock \emph{Neurocomputing}, 416:340--351, 2020.

\bibitem{ruptures}
C.~Truong, L.~Oudre, N.~Vayatis.
\newblock Selective review of offline change point detection methods.
\newblock \emph{Signal Processing}, 167:107299, 2020.

\bibitem{kliep}
M.~Sugiyama, S.~Nakajima, H.~Kashima, P.~von B\"unau, M.~Kawanabe.
\newblock Direct importance estimation with model selection and its
         application to covariate shift adaptation.
\newblock In \emph{NIPS}, 2008.

\bibitem{sugiyama_book}
M.~Sugiyama and M.~Kawanabe.
\newblock \emph{Machine Learning in Non-Stationary Environments: Introduction
         to Covariate Shift Adaptation}.
\newblock MIT Press, 2012.

\bibitem{rcssl_rules}
RoboCup Small Size League.
\newblock Rules of the RoboCup Small Size League. \\
\newblock \url{https://ssl.robocup.org/league-software/}

\end{thebibliography}

\end{document}
